\documentclass[10pt]{article}
\usepackage{amsmath,booktabs,hyperref,placeins}
\makeatletter
\def\input@path{{../shared/}}
\makeatother
\usepackage{hkustgz-paper}
\title{Model-Authored Improvers\\for Battery-Free Design}
\papersubtitle{Phase 4 local experimental report, 20 October 2026}
\author{Working draft}
\date{}
\begin{document}
\maketitle
\begin{abstract}
Can changing a method-proposal procedure improve battery-free design under a
fixed scientific task and a complete experimental cost boundary? We compare
ordinary reflection with one model-authored replacement using three paired
descendant lineages, a shared historical packet, and fresh design episodes.
All 30 effective selection cells pass. Among 27 final cells, the unchanged
base method, ordinary-reflection descendants, and replacement descendants
obtain 7/9, 6/9, and 8/9 complete-service PASS outcomes. All 18 cells on two
digital derivative tasks pass; differences occur only on the original event
service under a previously unexposed waveform. Four non-submissions end at a
32,768-token response limit and remain in denominators and costs. Full actual
ordinary-reflection and replacement pipelines cost 135 and 138 requests;
after removing only four retained apparatus incidents, they cost 123 and 129.
The 77 actual attempts use 367 requests in total. This is a descriptive result
conditional on one shared improver formation, not independent training
replication or a stable advantage. Default method promotion, recursive
self-improvement, and physical qualification are not established.
\end{abstract}
% Local phase 4 report. AUTO blocks use only pinned closed results.
\section{Experimental question and objects}

The experiment asks whether changing the procedure that proposes a design
method changes the service achieved by its descendants and the resources
needed to obtain them. Three objects have distinct roles. A design $D$ is a
task-specific device or timing artifact. A method $H$ is procedural text
supplied at design-task entry. An improver $G$ is the proposal procedure used
to generate $H$; it is neither a device submission nor a method already in
force. Writing a new $G$ establishes an authored artifact. Establishing a
useful change requires fresh design decisions, an unchanged service evaluator,
and costs that include formation, proposal, selection, and execution.

The reference $H_0$ and ordinary-reflection procedure $G_0$ are inherited
unchanged from phase 3. One prospective episode formed $G_1$. Its allowed
change was confined to the natural-language proposal procedure. Tasks,
physical-model semantics, default $H_0$, proposal prefix, source access, tools,
model configuration, budgets, evaluator, and selection rule remained fixed.
There was no weight update. Three descendant lineage pairs subsequently used
the one frozen $G_1$. They are repeated descendants conditional on a shared
formation, not three independent improver training events.

The registered finite experiment comprises 73 logical cells and, after an
explicit cache-representation amendment, 77 actual attempts. All are closed.
The reference method and the two descendant arms obtain 7/9, 6/9, and 8/9
final PASS outcomes, respectively. The numerical difference is confined to
one original-service task; it is not evidence of a stable improver advantage.
This report binds the protocol, amendment, original artifacts, selection
freeze, complete results, and actual accounting~\cite{protocol,results}.

\section{Information, pairing, and frozen decisions}

\subsection{Common historical information}

Both proposal arms receive the same 75,955-byte mechanically projected model
history. The complete 257,661-byte packet remains bound for host audit; it is
not the packet delivered to the model. This distinction matters when
describing available evidence. The histories contain earlier public
experience, including failures and qualification limits, rather than final
test feedback. The separate original-service development transfer, O0, is
excluded from $G_1$ formation and all $G/H$ proposal and selection inputs.

The configured provider reports the alias \texttt{deepseek-flash}. The actual
configuration is hash-bound, but independently pinned model weights remain
UNKNOWN. The claim is therefore conditional on the recorded configured
provider, not a verified immutable weight checkpoint. The public tool
projection is common to all arms: arrays longer than twelve expose their
count, full canonical hash, first six entries, and last six entries. Host
records retain the full response. Public feedback and host service evaluation
remain separate.

\subsection{Fresh descendants and design episodes}

For lineage $\ell\in\{1,2,3\}$ and arm $g\in\{0,1\}$, two named sibling
methods are proposed independently:
\[
 H_{\ell,g,j}\leftarrow
 \operatorname{Propose}(G_g,H_0,\mathcal{E}_{\mathrm{public}}),
 \qquad j\in\{1,2\}.
\]
Each sibling starts from $H_0$ and the identical history, rather than editing
its sibling's output. This gives twelve actual $H$ proposals. A consumed
descendant consists of unchanged $H_0$, two newlines, and the original
proposed method text. Its hash therefore differs from the proposal-text hash;
no host rewriting is implied. Each eligible method receives two fresh design episodes on
\texttt{b-pulse-fast} and \texttt{b-volatile-reset}. Six fresh $H_0$ control
episodes provide the same two tasks in each lineage. Thus selection comprises
twenty-four descendant cells and six common-control cells.

Every $D$ episode starts at task entry with the registered tools, information,
caps, and frozen initial exact-input cache, plus an isolated overlay. A cache
hit reuses a physical-model response for exactly matching input; it does not
reuse an Agent decision. It still consumes the registered logical design
exposure. A changed input cannot inherit the old physical response. Repeated
Agent decisions that share one cached response are not independent physical
experiments.

All thirty effective selection cells closed, and all six arm selections were
frozen before any final input was delivered to a participant. The sealed
preservation audit verifies this ordering and actual payload
derivation~\cite{selection,preservation}. Sibling ranking uses, in order, more complete-service PASS outcomes,
fewer wrong service judgments, fewer charged requests for the proposal plus
its two effective selection episodes, and stable child index. If neither
sibling is admissible, the protocol requires an explicitly labeled $H_0$
fallback. Invalid, no-update, failed, and censored attempts remain recorded.
They do not authorize replacement proposals or a change to the default $H_0$.

The final comparison contains twenty-seven fresh $D$ episodes:
three final tasks for six selected descendants and three lineage-matched
$H_0$ controls. These tasks cover a within-model battery-free derivative, a
C-lite service derivative, and the original service under a separately
frozen, previously unexposed waveform. Unexposed means absent from formation
and selection; a final $D$ participant normally receives the complete frozen
current-task facts, including its waveform. This is not a hidden-waveform
test of the final designer. Final outcomes were not used to reselect, tune,
repair, or promote a method.

\section{What the model actually formed}

\subsection{One frozen proposal procedure}

The actual $G_1$ is a 4,599-character procedure. It asks future proposals to
start with a source and capability ledger, express a typed symbolic method,
and branch on the consuming world's latch and source characteristics. It
then asks the proposer to triage recorded delivery, rehearse the procedure
against retained history, state a source-bound and tick-bound falsifier,
retain errors and costs, use the proposal channel appropriate to $H$, and
return no-update when support is absent. Historical numerical examples are
included as examples; they are not a new $D$ submission.

These are descriptions of the model's instructions, not accepted scientific
rules. In particular, its blanket treatment of events after the last
transmission as benign can fail for ongoing or multiple-event services.
Required rehearsal, exact source fields, and mandatory UNKNOWN judgments can
also constrain useful action or fail when a consuming interface differs.
The procedure and its descendants are preserved as submitted; the host does
not silently correct their formulas, channel assumptions, or service logic.

\subsection{Closed formation evidence}

One actual provider episode formed $G_1$ with three requests. Twelve actual
provider episodes then submitted twelve proposed methods, using forty-three
requests. Raw and normalized submissions are retained. Host readback binds
each descendant to $H_0$, its assigned $G$, the common model history, and its
submitted method hash. All twelve received artifact-scope semantic admission:
they change a procedural $H$ rather than submitting a device or modifying
the apparatus. Every audit explicitly leaves scientific validity
\texttt{NOT\_ESTABLISHED}. Admission does not certify the proposed arithmetic,
rehearsal claims, transfer, or eventual service performance.

% BEGIN AUTO formation
\begin{table}[htbp]
\centering
\small
\begin{tabular}{lrrr}
\toprule
Component & Episodes & Requests & Upper USD \\
\midrule
Shared $G_1$ formation & 1 & 3 & 0.0269157 \\
$G_0$ descendants & 6 & 20 & 0.1832424 \\
$G_1$ descendants & 6 & 23 & 0.2658975 \\
Formation and proposals & 13 & 46 & 0.4760556 \\
\bottomrule
\end{tabular}
\caption{Actual formation and proposal costs only; selection, final execution, controls and host incidents are additional.}
\label{tab:formation}
\end{table}
% END AUTO formation

The $G_0$ descendants contain 6,215--9,014 Unicode characters of method text;
the $G_1$ descendants contain 9,300--11,054. Longer text is an observed artifact
property in this sample, not evidence of a better method. Summed episode wall
time is 38.723294 seconds for formation, 352.501470 for the six $G_0$
proposals, and 510.444238 for the six $G_1$ proposals. These sums are not
elapsed campaign time.

Gateway logs contain fifteen \texttt{UNKNOWN\_PUBLIC\_SOURCE} failures across
five $H$ episodes, three in each. The failed acquisitions and their costs
remain in the original episodes. The thirteen formation/proposal ledgers have
no unknown or missing provider usage. A tool reporting an unavailable source
is distinct from UNKNOWN billing usage. Participant accounts of schema
rejection or native calls are not substituted for trusted logs: the retained
audit notes an unestablished schema-rejection claim and a conflation of source
acquisition with native engineering dispatch. Formation/proposal host records
report zero native dispatches.

\section{Apparatus incident and registered recovery}

An explicitly retained hold snapshot records three selection episodes with
\texttt{KeyError: 'trace'} and a fourth censored episode. The historical cache
loader supplied \texttt{observations.json} without the separately sealed
\texttt{trace.csv} required by assessment. This is a representation failure in
the apparatus; it cannot establish inferiority of either improver arm. At
that snapshot, nine selection attempts had closed, five were valid, and no
final episode had closed. These are historical snapshot counts, not a claim
about the current campaign state.

The approved amendment changes historical CSV loading and provenance
checks~\cite{amendment}.
It preserves original v1 records, compressed raw responses, task semantics,
methods, participant payloads, model, service, and the original tighter request
and wall-time ceilings. Six existing-data qualification checks passed: one
restored historical assessment, invariance of all five valid selections, and
rejection of missing trajectory, wrong trajectory binding, wrong exact input,
and wrong world. Those recorded checks used no new provider request or
simulation; they are not hardware qualification.

Four prospectively named \texttt{recovery-v11} identities then ran with fresh
contexts and the original payload and initial cache. Their sealed receipt
records four actual complete-service PASS outcomes, twenty provider requests,
and an upper cost of USD 0.1182666. This establishes the reported recovery
outcomes only. The four original faulty/censored attempts remain separately
visible: twenty-one requests and USD 0.1318473 of host-incident overhead. They
are not overwritten, relabeled as method failures, or erased from arm and
campaign cost. Neither $G_1$ nor any $H$ was regenerated. The amendment retains
seventy-three canonical logical slots while permitting at most seventy-seven
actual attempts, including the four original attempts.

\section{Selection cost and full pipeline cost}

Let $c(e)$ denote the recorded resource vector for episode $e$: charged
requests, summed episode wall time, and settled ledger upper USD. Formal
sibling ranking uses the proposal and its two effective selection episodes.
Under the amendment, each affected canonical selection slot resolves to its
named recovery identity; original incident attempts are separately itemized.
This formal ranking cost must not be mistaken for the cost of obtaining the
arm's result. Full arm accounting requires
\[
 C_{\ell,g}=A_{\ell,g}
 +\sum_{j=1}^{2}\!\left[c(P_{\ell,g,j})
                   +\sum_{s=1}^{2}c(S^{\mathrm{eff}}_{\ell,g,j,s})\right]
 +\sum_{t=1}^{3}c(F_{\ell,g,t})+c(I_{\ell,g}),
\]
where $A$ is allocated formation cost, $P$ a proposal, $S^{\mathrm{eff}}$
an effective selection episode, $F$ a final episode under the frozen selected
method or fallback, and $I$ the retained incident attempts. Both siblings,
including the unselected sibling, are charged. Common $H_0$ controls and the
separate initial transfer experiment are accounted once at campaign level,
not copied into every arm or omitted from a reported campaign total.

The single $G_1$ event is counted once globally. Its three requests allocate
one to each $G_1$ lineage; USD 0.0269157 allocates USD 0.0089719 each. Wall time
is divided in conserved decimal thirds, with the last allocation absorbing
the division residue. These allocations do not create additional formation
events. Each arm is capped at 120 requests, 6,600 summed wall seconds, and
USD 10, including allocated formation and failures. An $H$ proposal has eight
requests and 400 seconds; a $D$ episode has twelve requests, 800 seconds,
forty tool calls, and two distinct logical designs. Global ceilings remain
860 requests, 54,400 summed wall seconds, USD 100, and 150 new scientific
executions. Estimates and ledger upper bounds remain distinct from invoices;
human and orchestration assistance remains nonzero and unpriced.

\section{Complete results}

\subsection{Separate original-service development transfer}

The registered O0 development experiment contains a candidate and fixed
periodic and rich-state controls~\cite{O0}. The candidate and rich-state
control pass nominal complete event service, while the periodic control
fails. The candidate's own judgment is UNKNOWN despite host PASS
(Table~\ref{tab:O0}). Matching the rich-state control establishes no benefit
beyond it. These three episodes use eighteen requests and USD 0.1579020
upper; their results are excluded from the common $G/H$ history. Original
service contract, physical UNKNOWN, and earlier A/R failures are retained.

% BEGIN AUTO O0
\begin{table}[htbp]
\centering
\small
\begin{tabular}{lllrr}
\toprule
Condition & Host service & Participant & Requests & Upper USD \\
\midrule
candidate & PASS & UNKNOWN & 8 & 0.0760389 \\
fixed-periodic & FAIL & FAIL & 5 & 0.0321168 \\
rich-state & PASS & PASS & 5 & 0.0497463 \\
\bottomrule
\end{tabular}
\caption{The separate O0 development transfer. UNKNOWN disagrees with the candidate's host PASS under the nominal judgment metric. O0 contributes no history to $G/H$ formation or selection.}
\label{tab:O0}
\end{table}
% END AUTO O0

\subsection{Selection and final service}

All 30 effective selection cells pass complete service with zero nominal
judgment disagreements. Thus sibling ranking reduces to proposal-plus-two-
selection request count, then stable child index. Lineages 2 and 3 under
$G_0$ tie on requests and select child 1; USD is not a tie breaker. The other
four arms select child 2 (Table~\ref{tab:selection}). All twelve siblings
remain eligible; none is discarded or repaired for predicted final
performance. The six selected composite method hashes are frozen before
final exposure~\cite{selection,preservation}.

% BEGIN AUTO selection
\begin{table}[htbp]
\centering
\small
\begin{tabular}{lrrrl}
\toprule
Arm & $q(h_1)$ & $q(h_2)$ & Selected & Consumed $H$ hash \\
\midrule
L1-$G_0$ & 14 & 13 & 2 & \texttt{76060e3a5bb3} \\
L1-$G_1$ & 15 & 12 & 2 & \texttt{ba3c85390ea1} \\
L2-$G_0$ & 13 & 13 & 1 & \texttt{bc7d9af16574} \\
L2-$G_1$ & 14 & 13 & 2 & \texttt{206485c17bc5} \\
L3-$G_0$ & 15 & 15 & 1 & \texttt{dde0e5465cf7} \\
L3-$G_1$ & 14 & 13 & 2 & \texttt{1d206feb1c02} \\
\bottomrule
\end{tabular}
\caption{All twelve siblings pass both selection tasks, with zero wrong service judgments; all six $H_0$ controls also pass (30/30 effective cells). Requests include each child's proposal and two effective selections. Tied request counts use stable child index, not USD. Consumed $H$ is unchanged $H_0$ followed by the original proposed text; hashes refer to this composite.}
\label{tab:selection}
\end{table}
% END AUTO selection

The 27 final cells produce 21 PASS, two submitted FAIL, and four incomplete
non-submissions (Table~\ref{tab:outcomes}). The battery-free stretched-window
task and C-lite half-power query task each pass 9/9, totaling 18/18 across
arms. Only the original event service under its unexposed development/selection
waveform separates outcomes: $H_0$ yields one PASS, one submitted FAIL and one
censored cell; $G_0$ yields zero, one and two; $G_1$ yields two, zero and one.
Appendix~\ref{app:cells} lists every final cell's design parameters, requests
and cost. The host's complete-service judgment remains separate from the
participant's assessment~\cite{results}.

% BEGIN AUTO outcomes
\begin{table}[htbp]
\centering
\footnotesize
\begin{tabular}{lrrrrr}
\toprule
Arm & Cells & PASS & Submitted FAIL & Censored & Disagreements \\
\midrule
$H_0$ & 9 & 7 & 1 & 1 & 1 \\
$G_0$ & 9 & 6 & 1 & 2 & 1 \\
$G_1$ & 9 & 8 & 0 & 1 & 0 \\
\bottomrule
\end{tabular}
\caption{All 27 registered final cells, including four token-censored non-submissions. Disagreements compare a submitted participant judgment with the host service result; censored cells have no submitted judgment.}
\label{tab:outcomes}
\end{table}
% END AUTO outcomes

Both submitted O failures, L1-$H_0$ and L3-$G_0$, are judged PASS by their
participants but rejected by the host for \texttt{WRONG\_EVENT\_TIME}. Both
use one-second sampling. The two successful $G_1$ O designs use 0.75-second
sampling, and L2-$H_0$ also succeeds at 0.75 seconds. The observed design
choice is therefore not a mechanism unique to $G_1$. Common submitted O
choices include 15 mF storage, nonvolatile event cache, and a confirmed-events
program; retries differ. These digital results do not establish a new
memory mechanism or broad service transfer.

The four incomplete cells are L1-$G_0$, L2-$G_0$, L2-$G_1$ and L3-$H_0$ on
the O task. Each final assistant response has stop reason \texttt{length},
32,768 output tokens and 32,768 reasoning tokens, with no final submission.
Their actual result files confirm no model-submitted artifact. They are
response-token-censored cells, not host cache failures and not submitted
service FAIL outcomes. Each remains in the nine-cell arm denominator and in
all charged costs; no replacement was run. Dropping them would change the
registered comparison.

\subsection{Costs of obtaining the result}

Full actual $G_1$ and $G_0$ pipelines use 138 and 135 requests, with upper
estimates USD 1.2447618 and USD 1.2585825. The lower full USD for $G_1$ does
not establish an efficiency advantage: original host incidents cost $G_0$
twelve requests/USD 0.0742245 and $G_1$ nine/USD 0.0576228. Removing only those
four incidents, while retaining all recoveries, gives $G_1$ 129 requests/USD
1.1871390 and $G_0$ 123/USD 1.1843580. The direction of the USD difference
changes. Common $H_0$ selection and final controls use 76 requests/USD
0.5154765 (Table~\ref{tab:costs}). This comparison charges learned arms for
forming and selecting methods; it does not make $H_0$ information exposure
identical to their longer historical text.

% BEGIN AUTO costs
\begin{table}[htbp]
\centering
\footnotesize
\begin{tabular}{lrrrr}
\toprule
Scope & Full $q$ & Full upper USD & Effective $q$ & Effective upper USD \\
\midrule
$G_0$ & 135 & 1.2585825 & 123 & 1.1843580 \\
$G_1$ & 138 & 1.2447618 & 129 & 1.1871390 \\
$H_0$ controls & 76 & 0.5154765 & 76 & 0.5154765 \\
\bottomrule
\end{tabular}
\caption{Full arm pipelines include both siblings, proposal/selection/final calls, retained failures, and allocated $G_1$ formation. Effective pipelines remove only the four original host incidents, retaining all four charged recoveries. $H_0$ has selection/final controls but no proposal or formation stage.}
\label{tab:costs}
\end{table}
% END AUTO costs

Final episodes alone cost $H_0$, $G_0$, and $G_1$ 45, 40, and 45 requests,
respectively, with USD 0.3416751, 0.4800981, and 0.5159961 upper. These are
not full pipeline costs. Fewer calls in a censored cell also cannot be read
as greater task efficiency. Per-lineage full costs appear in
Table~\ref{tab:lineage-costs}; neither unselected siblings nor failed,
censored, or confirmation calls are omitted.

Across all stages, the campaign closes 77 attempts with 367 uniquely charged
provider requests, USD 3.1767228 upper, and 5,951.639484 summed episode wall
seconds (Table~\ref{tab:stages}). Full $G_0$ plus $G_1$ pipelines, $H_0$
controls, and separate O0 development conserve these campaign totals.
There are 25 new digital scientific responses, zero native dispatches, and
zero new UNKNOWN provider-usage records. S0's sealed preservation readback
also retains 24 earlier UNKNOWN records; zero new UNKNOWN does not erase
them~\cite{preservation}. Exact-input reuse remains distinct from fresh Agent
decisions. Codex assistance and historical acquisition are unpriced, so the
ledger is not an invoice or complete end-to-end cost.

% BEGIN AUTO stages
\begin{table}[htbp]
\centering
\footnotesize
\begin{tabular}{lrrrr}
\toprule
Stage & Attempts & Requests & Upper USD & Summed wall s \\
\midrule
Original-service development & 3 & 18 & 0.1579020 & 240.418 \\
$G_1$ formation & 1 & 3 & 0.0269157 & 38.723 \\
Twelve $H$ proposals & 12 & 43 & 0.4491399 & 862.946 \\
Selection, including incidents & 34 & 173 & 1.2049959 & 2027.954 \\
Final evaluation & 27 & 130 & 1.3377693 & 2781.598 \\
Total & 77 & 367 & 3.1767228 & 5951.639 \\
\bottomrule
\end{tabular}
\caption{Conserved campaign accounting. The 34 selection attempts comprise 30 effective cells and four original incidents. Summed episode wall time is not elapsed campaign time. Rates and ledger upper estimates are not invoices.}
\label{tab:stages}
\end{table}
% END AUTO stages

Original host-incident overhead appears only in lineage 1. Table~\ref{tab:lineage-costs}
separates full and effective costs in every arm; the three $G_1$ shares
conserve the one actual formation event.

% BEGIN AUTO lineage_costs
\begin{table}[htbp]
\centering
\footnotesize
\begin{tabular}{lrrrrr}
\toprule
Lineage arm & PASS & Full $q$ & Full USD & Effective $q$ & Effective USD \\
\midrule
L1-G0 & 2/3 & 51 & 0.4392936 & 39 & 0.3650691 \\
L1-G1 & 3/3 & 53 & 0.4813594 & 44 & 0.4237366 \\
L2-G0 & 2/3 & 40 & 0.4202280 & 40 & 0.4202280 \\
L2-G1 & 2/3 & 41 & 0.3678073 & 41 & 0.3678073 \\
L3-G0 & 2/3 & 44 & 0.3990609 & 44 & 0.3990609 \\
L3-G1 & 3/3 & 44 & 0.3955951 & 44 & 0.3955951 \\
\bottomrule
\end{tabular}
\caption{Per-lineage full and effective costs. Formation is one event with three allocated shares, not three independent formations. Full/effective definitions are identical to Table~\ref{tab:costs}.}
\label{tab:lineage-costs}
\end{table}
% END AUTO lineage_costs
\FloatBarrier

\section{Limitations and interpretation}

The $G_0$--$G_1$ comparison holds historical information common, but observes
only one formed $G_1$, three descendant pairs, one provider alias, and a small
fixed digital task family. It cannot estimate variation across independently
formed improvers or establish broad autonomous improvement. Comparing an
informed descendant $H$ with $H_0$ also changes information embedded in method
text. That contrast is not a pure procedural effect. Historical rehearsal
can overfit retained observations; task-specific source fields and
proposal-role instructions can fail in a new consuming role. The cache
incident and explicit recovery must remain part of the interpretation and
cost record.

Physical qualification remains UNKNOWN. Earlier fixed-space A/R service
\texttt{NO\_GO}, dark-service \texttt{FAIL}, strict FIFO \texttt{FAIL}, and product
\texttt{NOT\_ACCEPTED} remain retained evidence boundaries. Digital service
outcomes, source readback, structural admission, and a successful cache repair
do not establish hardware acceptance. The 8/9 $G_1$ total is descriptive,
conditional on one formed improver and this task sample. No statistical
superiority, stable advantage, general transfer capability, or recursive
self-improvement is established. The unchanged $H_0$ remains the default;
neither the global method nor the improver is promoted.

\section{Reproducibility and conclusion}

The binding pins the complete handoff, source results, request ledgers, six
selections, frozen methods, and unchanged style/vector logo. The reader
recomputes counts, conserved costs, ranking, consumed method composition and
table text without a model or scientific rerun. Canonical JSON digests remain
distinct from byte hashes; original preparation-error receipts are preserved.

Run \texttt{render\_phase4.py} with \texttt{check} or \texttt{build}, using
the existing runtime and \texttt{-B} from this paper directory.
The binding/handoff records absolute commands, fixed TeX dates, two-pass build,
fonts, PDF hash/pages and visual QA. S0's single registered final repository
check remains pending after import; no SDK, science, init or CI was rerun here.
This local report does not assess literature novelty or submission readiness.

The experiment produces one model-authored proposal procedure and twelve
actual descendants under controlled information and cost accounting.
Selection saturates at PASS, final task differences occur only on O, and
four token-censored attempts materially affect completion. The retained
apparatus incident changes the full/effective USD comparison. These facts
support preserving the default method and the bounded evidence record,
rather than declaring an autonomous improvement capability.

\renewcommand{\refname}{Local sealed source records}
\begin{thebibliography}{9}
\bibitem{protocol} BF campaign. Phase 4 finite protocol v1.1 and original
registration, 20 October 2026. Local records: \nolinkurl{protocol-v1.1.json},
\nolinkurl{registration.json}. Full byte hashes are in the paper binding.
\bibitem{results} BF campaign. Closed actual results, final details and
descriptive analysis, 20 October 2026. Local records:
\nolinkurl{actual-results.json}, \nolinkurl{final-details.json},
\nolinkurl{analysis-summary.json}; bound by \nolinkurl{paper-handoff.json}.
\bibitem{selection} BF campaign. Selection completion, selected methods and
final selection freeze, 20 October 2026. Local records:
\nolinkurl{selection-complete.json}, \nolinkurl{selected-methods.json},
\nolinkurl{final-selection-freeze.json}.
\bibitem{preservation} BF campaign. Final source, usage and preservation
readback, 20 October 2026. Local record: \nolinkurl{preservation-final.json}.
\bibitem{amendment} BF campaign. Host incident, cache-trace amendment,
qualification and four recovery outcomes, 20 October 2026. Local records:
\nolinkurl{host-error-hold.json}, \nolinkurl{amendment-registration.json},
\nolinkurl{cache-repair-qualification.json}, \nolinkurl{recovery-results.json}.
\bibitem{O0} BF campaign. Three original-service development transfer
episodes, 20 October 2026. Local record: \nolinkurl{O0-transfer-results.json}.
\end{thebibliography}

\clearpage
\appendix
\section{Every final design cell}\label{app:cells}

The task identifiers are:
\begin{itemize}
\item Task 1: \nolinkurl{p4-bf-stretched-port-window-v1}.
\item Task 2: \nolinkurl{p4-clite-half-power-query-v1}.
\item Task 3: \nolinkurl{p4-o-unexposed-waveform-v1}.
\end{itemize}
L identifies a descendant lineage slot,
conditional on the same shared $G_1$ formation. All registered cells, including
non-submissions, are shown. Parameters are rounded only for display; the
bound JSON retains exact submitted values.

% BEGIN AUTO final_cells
\begin{table}[htbp]
\centering
\footnotesize
\begin{tabular}{rrlclrr}
\toprule
Task & L & Arm & Result & Design parameters & $q$ & Upper USD \\
\midrule
1 & 1 & $H_0$ & P & 10 / 0.005 & 5 & 0.0221046 \\
1 & 1 & $G_0$ & P & 14 / 0 & 4 & 0.0378396 \\
1 & 1 & $G_1$ & P & 10 / 0 & 5 & 0.0352167 \\
1 & 2 & $H_0$ & P & 22 / 0 & 5 & 0.0289770 \\
1 & 2 & $G_0$ & P & 22 / 0 & 6 & 0.0786750 \\
1 & 2 & $G_1$ & P & 10 / 0 & 5 & 0.0256398 \\
1 & 3 & $H_0$ & P & 22 / 0.05625 & 5 & 0.0335634 \\
1 & 3 & $G_0$ & P & 22 / 0.05625 & 5 & 0.0416196 \\
1 & 3 & $G_1$ & P & 15 / 0.05625 & 5 & 0.0593052 \\
2 & 1 & $H_0$ & P & 40 / 6.5 & 5 & 0.0218922 \\
2 & 1 & $G_0$ & P & 40 / 7.5 & 5 & 0.0543483 \\
2 & 1 & $G_1$ & P & 50 / 6.5 & 5 & 0.0664434 \\
2 & 2 & $H_0$ & P & 50 / 6.5 & 5 & 0.0242127 \\
2 & 2 & $G_0$ & P & 50 / 6.89 & 5 & 0.0509238 \\
2 & 2 & $G_1$ & P & 50 / 6 & 5 & 0.0445269 \\
2 & 3 & $H_0$ & P & 50 / 6.5 & 5 & 0.0268155 \\
2 & 3 & $G_0$ & P & 50 / 7 & 4 & 0.0327336 \\
2 & 3 & $G_1$ & P & 50 / 6.5 & 5 & 0.0415404 \\
3 & 1 & $H_0$ & F & 1 / 1 & 5 & 0.0606894 \\
3 & 1 & $G_0$ & C & --- & 3 & 0.0549690 \\
3 & 1 & $G_1$ & P & 0.75 / 0.25 & 6 & 0.0990039 \\
3 & 2 & $H_0$ & P & 0.75 / 5 & 7 & 0.0753684 \\
3 & 2 & $G_0$ & C & --- & 3 & 0.0501078 \\
3 & 2 & $G_1$ & C & --- & 3 & 0.0581409 \\
3 & 3 & $H_0$ & C & --- & 3 & 0.0480519 \\
3 & 3 & $G_0$ & F & 1 / 0.25 & 5 & 0.0788814 \\
3 & 3 & $G_1$ & P & 0.75 / 1 & 6 & 0.0861789 \\
\bottomrule
\end{tabular}
\caption{Every final cell. P: submitted PASS; F: submitted FAIL; C: token-censored, no submission. Tasks 1/2 report capacitance ($\mu$F) / transmit gate (s); task 3 reports sample period (s) / retry (s). All submitted task-3 designs use 15 mF, nonvolatile cache and confirmed-events; other fields and complete JSON artifacts remain bound in final-details. Disagreements are task 3, L1-$H_0$ and L3-$G_0$, both judged PASS by the participant but FAIL by the host.}
\label{tab:final-cells}
\end{table}
% END AUTO final_cells
\FloatBarrier

\end{document}
